\section{Candidate Physical System and Experimental Test Strategy}
\label{sec:candidate_system}

The purpose of this section is to translate the general framework introduced above into a physically testable problem. Among known fissile systems, the low-lying isomeric state of $^{235}\mathrm{U}$, denoted $^{235m}\mathrm{U}$, is a particularly useful candidate because the two central ingredients of nuclear-state fission control are already partially established: the isomer can be produced and studied experimentally, and its neutron-induced fission response differs measurably from that of the nuclear ground state.

Importantly, the use of $^{235m}\mathrm{U}$ in connection with fission control is not itself a new suggestion. Previous work has explicitly noted that the isomer may be exploited to control nuclear fission because its thermal-neutron fission cross section differs substantially from that of ground-state $^{235}\mathrm{U}$ \cite{deer1988,liu2022}. The purpose of the present framework is therefore more specific: to treat the externally controllable nuclear-state population as a dynamical control variable, connect it explicitly to an effective fission cross section, and then propagate that microscopic modification into macroscopic quantities such as fission rate, neutron multiplication, and reactor reactivity.

\subsection{$^{235m}\mathrm{U}$ as a candidate system}

The first excited nuclear state of $^{235}\mathrm{U}$ lies only
\begin{equation}
E_m = 76.737 \pm 0.018\ \mathrm{eV}
\end{equation}
above the ground state \cite{ponce2018}. This exceptionally low excitation energy is unusual for a nuclear transition. The state has a half-life of approximately 26 min and decays predominantly through internal conversion.

For the present proposal, however, the most important property is not its low excitation energy but its different neutron-induced fission response.

D'Eer \textit{et al.} measured the fission cross section of $^{235m}\mathrm{U}$ relative to that of ground-state $^{235}\mathrm{U}$ and obtained
\begin{equation}
\frac{\sigma_f^m}{\sigma_f^g}=1.61\pm0.44
\end{equation}
for cold neutrons and
\begin{equation}
\frac{\sigma_f^m}{\sigma_f^g}=2.47\pm0.45
\end{equation}
for thermal neutrons \cite{deer1988}.

Thus, in this system,
\begin{equation}
\boxed{\sigma_f^m \neq \sigma_f^g}
\end{equation}
is an experimentally established fact.

This result is central because it demonstrates that the neutron-induced fission response of a fissile isotope is not determined by isotopic composition alone. It can depend strongly on the internal nuclear state.

The microscopic link
\begin{equation}
\rho_N \longrightarrow \sigma_f
\end{equation}
required by the framework developed in the previous sections therefore already exists experimentally for at least one relevant fissile system.

\subsection{Minimal two-population description}

Let $p_m(t)$ be the fraction of $^{235}\mathrm{U}$ nuclei occupying the isomeric state and
\begin{equation}
p_g(t)=1-p_m(t)
\end{equation}
the fraction occupying the ground state.

For a neutron of energy $E$, the effective microscopic fission cross section of the ensemble is
\begin{equation}
\sigma_f^{\mathrm{eff}}(E,t)
=
[1-p_m(t)]\,\sigma_f^g(E)
+
p_m(t)\,\sigma_f^m(E).
\end{equation}

Equivalently,
\begin{equation}
\sigma_f^{\mathrm{eff}}(E,t)
=
\sigma_f^g(E)
+
p_m(t)
\left[
\sigma_f^m(E)-\sigma_f^g(E)
\right].
\end{equation}

Introducing
\begin{equation}
r(E)=\frac{\sigma_f^m(E)}{\sigma_f^g(E)},
\end{equation}
the relative change with respect to a purely ground-state ensemble becomes
\begin{equation}
\frac{\Delta\sigma_f(E,t)}{\sigma_f^g(E)}
=
p_m(t)[r(E)-1].
\end{equation}

Using the measured thermal-neutron value $r\simeq2.47$,
\begin{equation}
\frac{\Delta\sigma_f}{\sigma_f^g}\simeq1.47\,p_m.
\end{equation}

This equation exposes the control problem directly. Once state-dependent fission cross sections are known, the only new dynamical quantity required at this level is the controllable isomeric fraction $p_m(t)$.

A state-control actuator would therefore act according to
\begin{equation}
u_Q(t)\longrightarrow p_m(t)\longrightarrow\sigma_f^{\mathrm{eff}}(t).
\end{equation}

The central microscopic feasibility question becomes whether $p_m(t)$ can be externally and reproducibly controlled, rather than whether the two states possess different fission probabilities, which is already experimentally established.

\subsection{External population of the isomer}

Producing a useful controlled population of $^{235m}\mathrm{U}$ remains the principal unresolved step.

Several nucleus-electron coupling mechanisms have been investigated for this unusually low-energy transition. Nuclear excitation by electronic transition (NEET) has received particular attention. Experimental claims have historically been inconsistent, and a dedicated 2016 search using a laser-generated uranium plasma observed no statistically significant $^{235m}\mathrm{U}$ production attributable to NEET. That work established an upper bound on the excitation rate under its experimental conditions \cite{chodash2016}.

A different route was investigated theoretically by Liu and Wang in 2022 using inelastic scattering of low-energy electrons. Their distorted-wave calculations predict excitation cross sections of order
\begin{equation}
10^{-38}-10^{-37}\ \mathrm{cm^2}
\end{equation}
for electron energies near the excitation threshold, several orders of magnitude larger than earlier plane-wave estimates \cite{liu2022}.

The same study noted that indirect population through higher nuclear levels may be more efficient than direct excitation of the $76.7$-eV transition. In particular, higher-lying levels coupled to the ground state by lower-multipolarity transitions provide possible intermediate pathways for population pumping \cite{liu2022}.

These results do not establish a practical actuator. They establish instead a well-defined open problem: whether a controllable excitation mechanism can produce a sufficiently large and sufficiently rapid modulation of $p_m$.

The proposed control framework does not require a particular excitation technology. The actuator may be electromagnetic, electron-mediated, plasma-mediated, or based on another externally tunable coupling mechanism. What matters is that it produces a reproducible mapping
\begin{equation}
u_Q(t)\rightarrow p_m(t).
\end{equation}

\subsection{Distinguishing excitation from control}

A crucial distinction must be maintained between \emph{producing} an isomer and \emph{controlling} the fission response through the isomeric population.

The existence of a process
\begin{equation}
^{235}\mathrm{U}_g\rightarrow{}^{235m}\mathrm{U}
\end{equation}
is not by itself sufficient to demonstrate nuclear-state fission control.

A control mechanism requires that an externally adjustable parameter $u_Q$ produce a predictable variation
\begin{equation}
p_m=p_m(u_Q),
\end{equation}
and that this variation subsequently produce a corresponding change in neutron-induced fission while the neutron field is otherwise held fixed.

The relevant chain is therefore
\begin{equation}
\boxed{
u_Q
\rightarrow
p_m
\rightarrow
\sigma_f^{\mathrm{eff}}
\rightarrow
R_f
}
\end{equation}
rather than simply
\begin{equation}
u_Q\rightarrow\text{nuclear excitation}.
\end{equation}

This distinction is what connects nuclear spectroscopy and nuclear excitation physics to control theory.

\subsection{Minimal falsification experiment}

The proposed mechanism can be tested without constructing a multiplying reactor system.

Consider two otherwise equivalent fissile ensembles exposed to the same neutron spectrum,
\begin{equation}
\Phi_1(E)=\Phi_2(E).
\end{equation}

Let the nuclear-state populations differ,
\begin{equation}
p_m^{(1)}\neq p_m^{(2)}.
\end{equation}

The corresponding fission-rate difference is
\begin{equation}
\Delta R_f
=
N_f
\int
\Phi(E)
\left[
p_m^{(2)}-p_m^{(1)}
\right]
\left[
\sigma_f^m(E)-\sigma_f^g(E)
\right]
\,dE.
\end{equation}

Since the state-dependent cross-section difference is already known to be non-zero for cold and thermal neutrons, the decisive experimental step is to generate the difference
\begin{equation}
\Delta p_m=p_m^{(2)}-p_m^{(1)}
\end{equation}
through an externally adjustable actuator.

The hypothesis predicts
\begin{equation}
\boxed{
\left.
\frac{\partial R_f}{\partial u_Q}
\right|_{\Phi}
\neq0
}
\end{equation}
provided that
\begin{equation}
\frac{\partial p_m}{\partial u_Q}\neq0
\end{equation}
and
\begin{equation}
\sigma_f^m-\sigma_f^g\neq0.
\end{equation}

This formulation also makes the experiment falsifiable. If the actuator produces a verified change in $p_m$ but no corresponding change in the neutron-induced fission response beyond experimental uncertainty, the proposed coupling model is incomplete or incorrect for that regime.

Conversely, observation of a reproducible correlation
\begin{equation}
u_Q\rightarrow p_m\rightarrow R_f
\end{equation}
with the neutron spectrum held fixed would constitute direct evidence of externally actuated nuclear-state control of fission.

\subsection{From microscopic control to macroscopic relevance}

A positive microscopic result would not, by itself, imply practical reactor control.

Let $f_Q$ denote the fraction of fissile nuclei whose state population can effectively be modified by the actuator. To first order, the macroscopic change in the fission response scales as
\begin{equation}
\frac{\delta\Sigma_f}{\Sigma_f}
\sim
f_Q\frac{\delta\sigma_f}{\sigma_f}.
\end{equation}

For the two-state $^{235}\mathrm{U}$ model, $f_Q$ may be absorbed into an effective ensemble-averaged isomeric fraction, but it is useful to retain the distinction when considering spatially nonuniform or incomplete actuation.

Macroscopic relevance therefore depends not only on the intrinsic difference
\begin{equation}
\Delta\sigma_f=\sigma_f^m-\sigma_f^g,
\end{equation}
but also on the achievable controlled population, the excitation and relaxation timescales, and the energetic cost of maintaining the desired state distribution.

A minimal feasibility parameter set is
\begin{equation}
\boxed{
\left\{
p_m^{\max},
\frac{dp_m}{dt},
\tau_m,
\Delta\sigma_f(E),
E_Q,
f_Q
\right\}.
}
\end{equation}

Only after these quantities are constrained experimentally does it become meaningful to propagate the effect into $\delta k_{\mathrm{eff}}$ and an effective reactor reactivity contribution $\rho_Q$.

\subsection{Scope of the present contribution}

The $^{235m}\mathrm{U}$ literature requires a careful statement of novelty.

Previous studies have already established that the isomer has a different neutron-induced fission cross section, and the possibility of exploiting this fact for fission control has been explicitly mentioned \cite{deer1988,liu2022}. The present work therefore does \emph{not} claim the first proposal to use $^{235m}\mathrm{U}$ to influence fission.

The distinct objective developed here is to formulate nuclear-state population as an explicit dynamical control variable and to embed it in a general multiscale chain,
\begin{equation}
u_Q(t)
\rightarrow
\rho_N(t)
\rightarrow
\sigma_f^{\mathrm{eff}}(E,t)
\rightarrow
R_f(t)
\rightarrow
k_{\mathrm{eff}}(t)
\rightarrow
\rho_Q(t).
\end{equation}

Within this framework, an isomer is not merely a nuclear state with unusual reaction properties. Its population becomes a candidate \emph{state variable of the controlled fission system}.

This formulation separates three questions that are often conflated:
\begin{enumerate}
\item \textbf{State dependence:} do different nuclear states have different fission responses?
\item \textbf{Actuation:} can their populations be externally and dynamically controlled?
\item \textbf{System coupling:} can the resulting microscopic modification influence macroscopic neutron multiplication?
\end{enumerate}

For $^{235m}\mathrm{U}$, the first question already has a positive experimental answer. The second remains an active nuclear-physics problem. The third is the reactor-physics and control-theory problem formulated in this work.

\subsection{Beyond $^{235m}\mathrm{U}$}

The framework is not restricted to a single isotope.

A suitable candidate system should satisfy as many as possible of the following conditions:
\begin{itemize}
\item at least two nuclear states or configurations with significantly different fission probabilities or fission cross sections;
\item an externally accessible transition or population-transfer pathway between those states;
\item a controlled-state lifetime sufficiently long compared with the actuation timescale;
\item an actuation mechanism capable of addressing a non-negligible fraction of the fissile ensemble;
\item a state-dependent response that survives ensemble averaging;
\item an energy cost compatible with producing a measurable net change in the fission dynamics.
\end{itemize}

Long-lived spectroscopic isomers, shape isomers, and selected excited configurations of actinides therefore constitute natural classes of systems for future investigation.

The search for candidate nuclei can be formulated as an optimization problem over two properties:
\begin{equation}
\boxed{
\text{controllability of the nuclear state}
\times
\text{sensitivity of fission to that state}.
}
\end{equation}

A system with a very large state-dependent fission effect but negligible controllability is unsuitable. Likewise, a readily controllable state is of little value if its fission response is indistinguishable from that of the uncontrolled ensemble.

The most promising systems are those in which both factors are simultaneously non-negligible.

\subsection{Experimental decision tree}

The proposed research program can therefore proceed through a sequence of increasingly stringent tests.

First, demonstrate or identify reproducible external control of a fission-distinct nuclear-state population,
\begin{equation}
u_Q\rightarrow\rho_N.
\end{equation}

Second, verify under a fixed neutron spectrum that the controlled state produces the predicted modification of neutron-induced fission,
\begin{equation}
\rho_N\rightarrow\sigma_f^{\mathrm{eff}}\rightarrow R_f.
\end{equation}

Third, determine whether the achievable controlled fraction and response magnitude are sufficient to produce a measurable perturbation of neutron multiplication,
\begin{equation}
R_f\rightarrow\delta k_{\mathrm{eff}}.
\end{equation}

Only the third stage requires coupling the microscopic nuclear-state dynamics to reactor kinetics.

This hierarchy deliberately separates the fundamental physics question from the engineering one. A negative result at the first or second stage would constrain or falsify the proposed control mechanism without requiring reactor-scale experiments. A positive result would provide the quantitative parameters needed to determine whether nuclear-state control can ever become relevant at macroscopic scale.
